\documentclass{ximera}

\input{../../preamble.tex}

\title{Rekenen met getallen}

\begin{document}

\begin{abstract}
In dit hoofdstuk herhaal je de rekenvaardigheden die je nodig hebt voor de rest van de cursus. Je leert nauwkeurig werken met gehele getallen, breuken, decimalen en percentages.
\end{abstract}

\maketitle

\section{Voorkennis}
\begin{itemize}
\item Optellen, aftrekken, vermenigvuldigen en delen
\item Kommagetallen
\item Breuken
\end{itemize}


\section{Leerdoelen}
\begin{itemize}
\item Rekenen met gehele getallen
\item Rekenen met breuken
\item Percentages toepassen
\item Afronden en schatten
\end{itemize}


\section{Waarom dit hoofdstuk?}

Een goede basis is noodzakelijk om later met functies, vergelijkingen en vectoren te kunnen werken.

\begin{example}
Bereken

\[
18+24-13
\]
\end{example}

\begin{exercise}
Bereken

\[
125+48=\answer{173}
\]
\end{exercise}

\end{document}